% © 2026 Ing. David Jaroš — CC BY-NC-ND 4.0
%
% This work is licensed under a Creative Commons Attribution-NonCommercial-NoDerivatives
% 4.0 International License (CC BY-NC-ND 4.0).
%
% License History: Earlier drafts (up to v0.3) were released under CC BY 4.0.
% From v0.4 onward, all material is released under CC BY-NC-ND 4.0 to protect
% the integrity of the theoretical work during ongoing academic development.

\documentclass[11pt,a4paper]{article}
\usepackage{amsmath, amssymb, amsthm}
\usepackage{hyperref}
\usepackage{geometry}
\geometry{margin=2.5cm}

\newtheorem{theorem}{Theorem}
\newtheorem{proposition}{Proposition}
\newtheorem{definition}{Definition}
\newtheorem{remark}{Remark}
\newtheorem{conjecture}{Conjecture}

\title{$B_{\mathrm{base}}$ from Spectral Determinants\\
\large Attempt E-2: Torus Laplacian, Zeta Function, and $N^{3/2}$ Scaling}
\author{Ing.\ David Jaroš}
\date{2026}

\begin{document}
\maketitle

\begin{abstract}
We investigate whether the coefficient $B_{\mathrm{base}} \approx 41.57$
in the UBT fine-structure constant formula arises as a spectral determinant
of the torus Laplacian acting on the compact $\psi$-directions.  Using the
spectral zeta function regularisation technique, we test whether the
$N_{\mathrm{eff}}^{3/2}$ scaling emerges from the functional determinant
$\det(-\Delta_{T^d})$.  This is attempt E-2 in the $B_{\mathrm{base}}$
programme.  All claims are labelled with their proof status.
\end{abstract}

\tableofcontents

% ======================================================================
\section{Spectral Determinants and Zeta Function Regularisation}
\label{sec:zeta}
% ======================================================================

\subsection{Spectral zeta function}

For a positive operator $A$ with eigenvalues $\{\lambda_n\}_{n \geq 1}$, the
spectral zeta function is:
\begin{equation}
\label{eq:spec_zeta}
  \zeta_A(s) = \sum_{n=1}^\infty \lambda_n^{-s}, \qquad \Re(s) > d/2,
\end{equation}
and the functional determinant is defined via analytic continuation:
\begin{equation}
\label{eq:log_det}
  \ln\det A = -\zeta_A'(0).
\end{equation}

\subsection{Torus Laplacian spectrum}

The Laplacian on a $d$-dimensional rectangular torus
$\mathbb{T}^d = \prod_{i=1}^d S^1_{R_i}$ has eigenvalues:
\begin{equation}
\label{eq:torus_spec}
  \lambda_{\mathbf{n}} = \sum_{i=1}^d \frac{n_i^2}{R_i^2},
  \qquad \mathbf{n} = (n_1,\ldots,n_d) \in \mathbb{Z}^d.
\end{equation}

For the isotropic case $R_i = R_\psi$ for all $i$:
\begin{equation}
  \lambda_{\mathbf{n}} = \frac{|\mathbf{n}|^2}{R_\psi^2}.
\end{equation}

% ======================================================================
\section{Spectral Zeta Function of the Torus}
\label{sec:torus_zeta}
% ======================================================================

\subsection{Epstein zeta function}

For the $d$-torus (excluding zero mode), the spectral zeta function
is the Epstein zeta function:
\begin{equation}
\label{eq:epstein}
  Z_d(s) = \sum_{\mathbf{n} \in \mathbb{Z}^d \setminus \{0\}}
  \frac{R_\psi^{2s}}{|\mathbf{n}|^{2s}}.
\end{equation}

For $d = 1$ (single $S^1_\psi$):
\begin{equation}
  Z_1(s) = 2 R_\psi^{2s} \zeta(2s),
\end{equation}
where $\zeta$ is the Riemann zeta function.

\subsection{Functional determinant of $S^1$}

Using $\ln\det(-\Delta_{S^1}) = -Z_1'(0)$:
\begin{equation}
  Z_1'(0) = 4 R_\psi^0 \zeta(0)\ln R_\psi + 2 R_\psi^0 \zeta'(0)
  = -2\ln R_\psi + 2\zeta'(0)
  = -2\ln R_\psi - \ln(2\pi).
\end{equation}

\begin{proposition}[$S^1$ functional determinant]\label{prop:s1_det}
  \begin{equation}
    \det(-\Delta_{S^1}) = 2\pi R_\psi.
  \end{equation}
  \textbf{Status: [L0]} — standard result.
\end{proposition}

\subsection{Functional determinant of $T^d$}

For the $d$-torus:
\begin{equation}
\label{eq:td_det}
  \ln\det(-\Delta_{T^d})' = -Z_d'(0).
\end{equation}

At $d = 3$ (imaginary-time torus $S^1_\psi \times S^1_\chi \times S^1_\xi$):
\begin{equation}
\label{eq:z3_s0}
  Z_3'(0) = \sum_{\mathbf{n}\neq 0} \ln |\mathbf{n}|^2 \cdot R_\psi^0
  = 2 \sum_{\mathbf{n}\neq 0} \ln|\mathbf{n}|.
\end{equation}

The sum $\sum_{\mathbf{n}\neq 0}\ln|\mathbf{n}|$ requires zeta-function
regularisation.  Using the Kronecker limit formula at $d=3$, one obtains
a finite result depending on $R_\psi$.

% ======================================================================
\section{Connection to \texorpdfstring{$N_{\mathrm{eff}}^{3/2}$}{Neff32}}
\label{sec:neff}
% ======================================================================

\subsection{Effective mode count from spectral determinant}

Define the effective mode count as the exponential of the log-determinant:
\begin{equation}
\label{eq:neff_def}
  N_{\mathrm{eff}} \;\stackrel{?}{=}\; [\det(-\Delta_{T^3})]^{1/3}.
\end{equation}

The $\tfrac{1}{3}$ power is chosen so that $N_{\mathrm{eff}}$ scales linearly
with $R_\psi$ for a single circle.

\begin{remark}[Scaling check]
  From Proposition~\ref{prop:s1_det}, $\det(-\Delta_{S^1}) = 2\pi R_\psi$.
  For a 3-torus, $\det(-\Delta_{T^3}) = (2\pi R_\psi)^3$ at leading order
  (independent dimensions).  Hence:
  \begin{equation}
    N_{\mathrm{eff}} = (2\pi R_\psi)^3)^{1/3} = 2\pi R_\psi.
  \end{equation}
  This gives $B_{\mathrm{base}} = N_{\mathrm{eff}}^{3/2} = (2\pi R_\psi)^{3/2}$.
  
  Numerically, $B_{\mathrm{base}} = 41.57$ would require:
  $(2\pi R_\psi)^{3/2} = 41.57 \Rightarrow R_\psi \approx 3.2 \cdot 10^{-9}$
  in natural units.  This is a prediction that can be checked against other
  constraints on $R_\psi$.
\end{remark}

\begin{remark}[Status]
  The spectral determinant approach gives a \emph{consistent} picture:
  the $3/2$ exponent arises naturally from the combination of (i) a
  $3$-torus determinant and (ii) the $1/3$ normalisation in
  \eqref{eq:neff_def}.  However, the normalisation $1/3$ is not derived
  — it is chosen to match dimensions.

  Status: \textbf{[L2]} — structurally consistent, not a proof.
\end{remark}

% ======================================================================
\section{Casimir Energy and One-Loop Determinant}
\label{sec:casimir}
% ======================================================================

\subsection{One-loop Casimir energy}

The one-loop effective potential for the Θ-field on $T^3$ contains the
Casimir energy:
\begin{equation}
\label{eq:casimir}
  V_{\mathrm{Casimir}} = -\frac{1}{2}\ln\det(-\partial^2 + m^2)
  = -\frac{1}{2} Z_{T^3}'(0)|_{m^2}.
\end{equation}

For $m = 0$ (massless limit) and $R_\psi \to 0$:
\begin{equation}
  V_{\mathrm{Casimir}} \sim -\frac{3\zeta(4)}{(2\pi)^4 R_\psi^4}
  = -\frac{\pi^2}{90 R_\psi^4}.
\end{equation}

\begin{remark}[Connection to one-loop $\alpha$]
  The Casimir energy contributes to the one-loop running of $\alpha$ through
  the vacuum energy density.  This is related to the $B_\alpha$ coefficient
  rather than $B_{\mathrm{base}}$.  See
  \texttt{consolidation\_project/appendix\_ALPHA\_one\_loop\_biquat.tex} §B.
  Status: \textbf{[L1]}.
\end{remark}

% ======================================================================
\section{Modular Transformation and Poisson Duality}
\label{sec:poisson}
% ======================================================================

Under the modular transformation $R_\psi \to 1/R_\psi$, the torus zeta
function satisfies Poisson duality:
\begin{equation}
\label{eq:poisson}
  Z_d(s; R_\psi) = R_\psi^{2s - d}\, Z_d(d/2 - s;\, 1/R_\psi).
\end{equation}

For $d = 1$ this is Riemann's functional equation for $\zeta(s)$.
For $d = 3$, \eqref{eq:poisson} relates large-$R_\psi$ to small-$R_\psi$:
\begin{equation}
  W(R_\psi) = R_\psi^{3/2}\, S(1/R_\psi),
\end{equation}
where $W$ is the winding-mode sum and $S$ is the spectral-mode sum.

\begin{remark}[Exponent $3/2$ from Poisson duality]
  The prefactor $R_\psi^{3/2}$ in the Poisson duality relation for $d=3$
  exactly produces the $3/2$ exponent.  This suggests that
  $B_{\mathrm{base}} \propto N_{\mathrm{eff}}^{3/2}$ might be a direct
  consequence of Poisson duality on the 3-torus.

  Source: \texttt{ubt\_with\_chronofactor/scripts/spectral/poisson\_duality\_demo.py}
  (verified: the prefactor is $+d/2$ for a Gaussian $f(x) = e^{-\pi|x|^2/\tau}$).

  Status: \textbf{[L2]} — structurally compelling but not yet a proof.
  The missing step is identifying $B_{\mathrm{base}}$ with the Poisson dual
  prefactor in the one-loop QED calculation.
\end{remark}

% ======================================================================
\section{Assessment and Next Steps}
\label{sec:assess}
% ======================================================================

\begin{center}
\begin{tabular}{ll}
\hline
Result & Status \\
\hline
$\det(-\Delta_{S^1}) = 2\pi R_\psi$ & \textbf{[L0]} proved \\
$\det(-\Delta_{T^3}) \sim (2\pi R_\psi)^3$ & \textbf{[L1]} leading order \\
$N_{\mathrm{eff}} = [\det]^{1/3}$ definition & \textbf{[L2]} motivated \\
$3/2$ from Poisson duality on $T^3$ & \textbf{[L2]} structurally compelling \\
Numerical coefficient of $B_{\mathrm{base}}$ from $\det$ & \textbf{[O]} open \\
\hline
\end{tabular}
\end{center}

The most promising direction from this investigation is the Poisson duality
observation (§\ref{sec:poisson}).  The next step is to compute the full
one-loop determinant with the KK sum and verify that the Poisson prefactor
$R_\psi^{3/2}$ appears in the coefficient of the log-running coupling, and
that this coefficient equals $B_{\mathrm{base}}$ numerically.

Continued in \texttt{research/B\_base\_heat\_kernel.tex}.

\end{document}
